% General theory: partially identified global portraits and downstream decisions.

The mapping calculus above can be stated more generally. Scientific sources
rarely reveal a unique global state. They reveal projections of that state,
with different coordinates observed, suppressed, coarsened, or measured under
different operationalizations. The appropriate integration object is
therefore not always one completed graph. It is an envelope of the global
portraits still compatible with the observations and their source-bound
constraints. This section develops that object and separates three questions:
what the sources identify, what additional information could identify, and
what a downstream decision may safely do before identification is complete.

The construction imports the logic of partial identification
\cite{manski2003partial,tamer2010partial}, comparison of information
\cite{blackwell1953equivalent}, and robust decision theory
\cite{berger1985decision,walley1991imprecise}. The contribution is not a new
claim to those donor theories. It is their claim-relative composition with
provenance-bound scientific projections, typed mapping constraints, and
global-portrait construction.

\subsection{Worlds, observation fibres, and portrait envelopes}

Let $\Omega$ be a set of scientifically admissible worlds. A world
$\omega\in\Omega$ fixes the source records, the source-to-projection
relations, and a global portrait $G(\omega)\in\mathcal G$. The observation
operator
\begin{equation}
  O:\Omega\longrightarrow\mathcal Y
\end{equation}
returns exactly the evidence available to the integration procedure: the
observed projection coordinates, explicit missingness, whatever source
identity/version information the interface exposes, and the mapping
constraints licensed by those sources. It does not fill an unobserved
coordinate from the world that generated it.

The operator is interface-relative. If a procedure receives source text and
semantic provenance, those are part of $O$; if it receives only projections,
they are not. Opaque record identifiers that carry no licensed scientific
meaning are quotiented by bijective relabelling rather than treated as a
lookup key for memorising gold. Every impossibility statement must name this
observation boundary.

\begin{definition}[Observation fibre and epistemic portrait envelope]
For observed evidence $y\in\mathcal Y$, define
\begin{equation}
  \Omega_y=\{\omega\in\Omega:O(\omega)=y\},\qquad
  \mathfrak E(y)=\{G(\omega):\omega\in\Omega_y\}.
\end{equation}
$\Omega_y$ is the observation fibre and $\mathfrak E(y)$ is the epistemic
portrait envelope. A query $q:\mathcal G\to\Theta$ has identified set
\begin{equation}
  \Theta_q(y)=\{q(g):g\in\mathfrak E(y)\}.
\end{equation}
The query is point identified at $y$ exactly when $\Theta_q(y)$ is a
singleton.
\end{definition}

The query can ask whether two projections may be glued, which typed
obstruction holds, what numerical estimand a merged record implies, or which
downstream action is warranted. Identification is therefore claim-relative:
an envelope can identify one query while leaving another unresolved. A
portrait need not be completed in every coordinate before it licenses a
specific conclusion.

The envelope also separates scientific plurality from epistemic uncertainty.
Plurality is a feature represented \emph{inside} a feasible portrait: several
incompatible but source-valid views coexist. Epistemic uncertainty is
variation \emph{across} feasible portraits because the available observation
does not decide which completion obtains. A plural-view terminal can therefore
be point identified; an unresolved terminal cannot.

\subsection{The fibre criterion and an impossibility result}

\begin{theorem}[Fibre criterion for scientific identification]
\label{thm:fibre}
For any observation $y$ with $\Omega_y\neq\varnothing$ and query $q$, the
following are equivalent:
\begin{enumerate}
  \item $q$ is point identified at $y$;
  \item $q\circ G$ is constant on $\Omega_y$;
  \item there exists an observation-measurable answer $d(y)$ that equals
        $q(G(\omega))$ for every $\omega\in\Omega_y$.
\end{enumerate}
If two worlds in the same observation fibre have different query values, no
deterministic procedure reading only $y$ can be correct in both worlds.
\end{theorem}

\begin{proof}
The image of $\Omega_y$ under $q\circ G$ is exactly $\Theta_q(y)$. It is a
singleton if and only if the map is constant on the fibre. In that case its
unique value defines $d(y)$. Conversely, if $d(y)$ is correct for every world
in the fibre, all those worlds must have query value $d(y)$. The final claim
follows immediately: the procedure receives the same input in both worlds and
therefore emits the same answer, which cannot equal two different values.
\end{proof}

This theorem turns a partial-observation failure into a scientific result
rather than an invitation to invent a fifth rule over the same information.
When a coordinate erased by extraction is the only fact distinguishing two
source relations, model capacity, prompt changes, and a more elaborate
comparison rule cannot restore it. The remedy must either acquire information,
declare an assumption that removes worlds from the fibre, or return a
set-valued/unresolved answer.

\subsection{A canonical sufficient interface}

The envelope is not a demand that every integration system share one internal
architecture. It defines the least claim-relative interface a downstream
consumer must be able to recover.

\begin{definition}[Envelope sufficiency]
A summary $S:\mathcal Y\to\mathcal Z$ is sufficient for query $q$ when there
exists a decoder $h$ such that
\begin{equation}
  \Theta_q(y)=h(S(y))\qquad\text{for every }y\in\mathcal Y.
\end{equation}
Thus equal summaries may never hide different identified sets.
\end{definition}

\begin{theorem}[Universal factorisation of sufficient interfaces]
\label{thm:factorisation}
The map $T_q:y\mapsto\Theta_q(y)$ is sufficient for $q$. Every other sufficient
summary $S$ refines it: there is a map $h$ with $T_q=h\circ S$. Consequently,
all sufficient implementations presented with the same downstream loss have
the same robust loss table and the same set of minimax actions, up to ties.
\end{theorem}

\begin{proof}
$T_q$ is decoded by the identity map. The factorisation for another sufficient
$S$ is its defining decoder. Robust upper loss for action $a$ is a function
only of the recovered set,
$\sup_{\theta\in T_q(y)}L(a,\theta)$. Equal recovered sets therefore give the
same value for every action, and minimisation gives the same optimal-action
set.
\end{proof}

This theorem explains both the ambition and the limit of the architecture. A
semantic product that exposes exactly the same identified set is not an
inferior baseline; it is an information-equivalent implementation and must
tie at this interface. A product that collapses two observations with
different identified sets is insufficient for that query, no matter how rich
its remaining representation is. Any empirical advantage must therefore come
from a genuinely finer observation, justified constraints, a more faithful
approximation to the envelope, or a different declared downstream loss---not
from the name or central location of the layer.

\subsection{Information refinement and the observation floor}

Say that $O_2$ refines $O_1$ when there is a map $r$ such that
$O_1=r\circ O_2$. Thus $O_2$ retains at least as much information as $O_1$;
the relation is deterministic and makes no probabilistic sufficiency claim.
Let $\mathcal A$ be a finite downstream action set and
$L:\mathcal A\times\Theta\to[0,\infty)$ a declared loss. Define the robust
decision floor
\begin{equation}
  B_q(y)=\min_{a\in\mathcal A}\ \sup_{\theta\in\Theta_q(y)}L(a,\theta).
  \label{eq:decision-floor}
\end{equation}

\begin{theorem}[Information-refinement monotonicity]
\label{thm:refinement}
If $O_2$ refines $O_1$, then for every realised world $\omega$,
\begin{equation}
 \Theta^{(2)}_q(O_2(\omega))
 \subseteq
 \Theta^{(1)}_q(O_1(\omega))
 \quad\text{and}\quad
 B^{(2)}_q(O_2(\omega))\leq B^{(1)}_q(O_1(\omega)).
\end{equation}
\end{theorem}

\begin{proof}
Every world having the refined observation $O_2(\omega)$ also has coarse
observation $r(O_2(\omega))=O_1(\omega)$, so the refined fibre is a subset of
the coarse fibre. Taking images under $q\circ G$ preserves inclusion. For any
fixed action, the supremum of its loss over a subset cannot increase; taking
the minimum over actions preserves the inequality.
\end{proof}

The quantity $B_q(y)$ is an observation floor, not a defect score for a
particular algorithm. It can fall when truthful extraction, measurement,
provenance, or source acquisition refines the observation. It cannot be made
to disappear merely by changing the decision rule over the same fibre. For a
candidate rule $\delta$, its floor-adjusted avoidable harm is
\begin{equation}
 X_\delta(y)=
 \sup_{\theta\in\Theta_q(y)}L(\delta(y),\theta)-B_q(y)\geq 0.
 \label{eq:excess-harm}
\end{equation}
This is the quantity a successor experiment should compare. Counting total
error without subtracting or stratifying by the observation floor confounds
avoidable decision failure with information that no compared rule received.

\subsection{Closed-form merge, split, and unresolved bounds}

For a binary mapping query, let $\theta=1$ mean that glue is licensed and
$\theta=0$ mean that it is not. The downstream actions are merge $M$, keep
separate $S$, and unresolved $U$. A false merge costs $c_{\mathrm{FM}}>0$, a
false split costs $c_{\mathrm{FS}}>0$, and deferral costs $c_U\geq0$.
Suppose an independently justified credal description gives only
\begin{equation}
  \Pr(\theta=1\mid y)\in[\underline p(y),\overline p(y)].
\end{equation}
No point probability is manufactured when this interval is wide.

\begin{proposition}[Robust downstream decision bound]
\label{prop:binary-bound}
The worst-case expected losses of the three actions are
\begin{equation}
 \overline R(M\mid y)=c_{\mathrm{FM}}(1-\underline p),\quad
 \overline R(S\mid y)=c_{\mathrm{FS}}\overline p,\quad
 \overline R(U\mid y)=c_U.
 \label{eq:binary-risk}
\end{equation}
Consequently, unresolved is minimax-optimal whenever
\begin{equation}
 c_U\leq
 \min\{c_{\mathrm{FM}}(1-\underline p),
       c_{\mathrm{FS}}\overline p\}.
 \label{eq:defer-bound}
\end{equation}
With no probabilistic authority beyond the vacuous interval $[0,1]$, this
reduces to
\[
 c_U\leq\min\{c_{\mathrm{FM}},c_{\mathrm{FS}}\}.
\]
\end{proposition}

\begin{proof}
Merging incurs loss only when $\theta=0$, so its expected loss is
$c_{\mathrm{FM}}(1-p)$ and is maximised at $p=\underline p$. Keeping separate
incurs loss only when $\theta=1$, so its expected loss is
$c_{\mathrm{FS}}p$ and is maximised at $p=\overline p$. Deferral has constant
loss. Comparing the three upper risks gives the criterion.
\end{proof}

The costs are downstream and use-specific. A clinical synthesis may assign a
larger cost to a false merge than a discovery interface; a screening system
may assign a larger cost to a false split. The portrait mechanism should
expose the identified set and provenance needed to evaluate those choices,
not silently encode one universal utility function.

\subsection{Global composition is a joint-completion problem}

Let each accepted source or mapping record impose a set
$C_j\subseteq\Omega$ of worlds consistent with its licensed conditions. The
jointly feasible set is
\begin{equation}
  F(y)=O^{-1}(y)\cap\bigcap_{j=1}^{m}C_j.
  \label{eq:joint-completion}
\end{equation}
A global glue requires $F(y)\neq\varnothing$; a claim additionally requires
its query to be constant on $F(y)$. Pairwise mapping success does not imply
either condition.

\begin{proposition}[Pairwise compatibility is insufficient]
\label{prop:pairwise}
There exist three mapping constraints that are pairwise jointly satisfiable
but have no common global completion. Hence a system that accepts every
pairwise-compatible edge can license an inconsistent global portrait.
\end{proposition}

\begin{proof}
Take $\Omega=\{1,2,3\}$ and
$C_1=\{1,2\}$, $C_2=\{2,3\}$, $C_3=\{1,3\}$. Every pairwise intersection is
nonempty, but $C_1\cap C_2\cap C_3=\varnothing$.
\end{proof}

Constraint processing and acyclic join theory provide stronger algorithms for
deciding global consistency than this elementary counterexample
\cite{dechter2003constraint}. The point here is architectural: pairwise
ontology or schema correspondences are proposals for $C_j$, not self-executing
authority for a global merge. Empty joint completion is an identified global
obstruction. Multiple completions with different query values are partial
identification. Multiple source-valid views inside every completion are an
identified plural portrait.

\subsection{Evaluation implied by the theory}

The envelope theory changes the scientific target from forced-label accuracy
to four jointly reported quantities:
\begin{enumerate}
  \item \emph{validity}: whether protected source-grounded truth is excluded
        from the returned envelope;
  \item \emph{sharpness}: the size or decision diameter of the envelope,
        reported only together with validity;
  \item \emph{floor-adjusted decision harm}: Equation~\eqref{eq:excess-harm}
        under a declared downstream loss; and
  \item \emph{information value}: the reduction in the identified set or
        decision floor from a specified observation refinement.
\end{enumerate}
These targets envelop the original glue-or-obstruction calculus. A singleton
envelope recovers an ordinary determinate mapping; a non-singleton envelope
explains why unresolved is required; an empty joint-completion set certifies a
global obstruction. The current experiments cover only special cases of this
theory. They do not by themselves establish calibrated naturalistic envelopes
or downstream utility.

\endinput
